\documentclass[11pt]{article}
\usepackage[T1]{fontenc}\usepackage{amsmath,amssymb,amsthm,lmodern,geometry}
\geometry{margin=.95in}\setlength{\emergencystretch}{2em}
\newtheorem{theorem}{Theorem}\newtheorem{lemma}{Lemma}
\newcommand{\dd}{\,d}\newcommand{\E}{\mathbb E}
\title{Inactive ancestry through a logarithmic resident interval}
\author{Analytic increment for review}\date{Version 1, November 2026}
\begin{document}\maketitle
\begin{abstract}
For the unclamped critical right-half entry in the saturated leading
system, a fixed set of initially inactive labels of mass at least
\(c_{b,\rho,\zeta}/(1+T)\) stays inactive throughout \([0,T]\).
The set depends on the horizon. The proof averages the per-rank ceiling
against the actual initial tail before constructing an invariant region.
At the opposite end of the distribution the Q4 atom initially moves
down in both coordinates, and crosses zero in finite time for an
explicit range of large scales. We give the initial return and weak
gate thresholds, and a two-coordinate weak-coupled ceiling with the
signed mean forcing included. Persistence of a nonvanishing fraction
on arbitrarily long windows remains open.
\end{abstract}

\section{State and source ranges}\label{an04impv1:scope}
Use \(13/20\le\rho\le7/10,\ \lambda=\rho^2,\ \ell=\log(1/q),\
\omega=(1-\lambda)/2,\ S_0=cq^{1/\omega},\ b=(c/A)^\omega>0\).
The offset \(\zeta>a\) stays in a compact subset of that domain;
\(a=\rho K(\rho a)\), \(K(t)=\phi(t)+t\Phi(t)\).
CFS laws and RPS q4 are positive-noise \emph{comparison} results
at \(T_A=\log(A/S_0)\). Their right-half limiting shape is now
prescribed as aligned leading data. Write
\[
 \langle j\rangle_b=\tfrac12j(a)+\frac2\pi\int_0^\infty
         j(H^{-1}(bu))(1+u^2)^{-2}\dd u,
\]
where \(H\) is CFS's increasing time-of-flight map from \((a,\infty)\)
onto \((0,\infty)\). This law has unit mass, with Q4 atom \(1/2\)
and continuous Q1 part of density asymptotic to \(2b^3/(\pi x^4)\).
Let \(m=\langle x\rangle_b\), and denote its decreasing upper-tail
quantile by \(\xi(f)\), \(0<f<1\); the order of the physical coordinates
is increasing in rank \(\sigma=1-f\). Endpoints of rank intervals
have zero rank measure. Equality at the probe belongs to the inactive set.

WIT popsystem/P leading concern a given positive leading measure with
sufficient moments, at finite times of characteristic evolution. We use
\(M=1,N=0\) until Section~\ref{an04impv1:weak}. RPS moments and
GOE ceiling concern the same field; the ceiling uses \(Y(0)\ge0\)
and the upper donor fraction. CFS spread uses co-order and common
radial reaction, giving \(y_j-y_i\ge y_j(0)-y_i(0)\) at \(M=1,N=0\).
The GOE initial gate threshold applies only to aligned input at the
stated mass; GOE gap excludes fixed positive mass with a sustained
positive exponential input lower rate, not fixed mass above a fixed gate.
No WIT learned-window or original-flow sign theorem is used here.

With \(t=\rho x,\ F(t,r)=\E[(Z+r)_+1_{Z>-t}]\), let
\(S(x)=\int F(\rho x,\rho\tilde x)d\mu,\ K_s=\int K(\rho x)d\mu,\
Y=\int y\,d\mu,\ X_m=\int x\,d\mu\).
At saturation the exact velocities are
\[
 2U=\rho\phi(t)-\rho S(x)+\lambda y\Phi(t),\qquad
 2V_s=\rho K(t)-Y,\qquad 2Y'=-Y+\rho K_s.
\]
In particular \(Y>0\) for finite time. Strong order follows by taking
differences in the same donor field; both off-diagonals in the
receiver Jacobian are \(\lambda\Phi(t)/2>0\).
For clarity the finite-interval estimates below apply also to the
unclamped law: \(D=1+\int(|x|+|y|)d\mu\) obeys \(D^+D\le7D\)
by \(S\le K_s\le\phi(0)+\rho\int x_+d\mu\) and
\(|U|+|V_s|\le3(|x|+|y|+D)\).
One may construct an evolution by clamping the initial rank map at
successively larger radii. These estimates and pointwise Gronwall give
\(|x|+|y|\le C_T(1+\xi(f))\) uniformly in the clamp. On each closed
rank interval bounded away from \(f=0\) the velocities and their receiver
Lipschitz constants are uniformly bounded (use
\(F_t(t,r)=(r-t)_+\phi(t)\)). Dependence on initial rank is therefore
equicontinuous there by the continuous initial quantile and Gronwall in
each common field. Compactness, a diagonal subsequence in rank and time,
and domination by \(C_T(1+\xi(f))\) in the donor integrals give the
unclamped characteristic equations on finite intervals. The proofs
apply to any such evolution, without comparing velocities across
different self-consistent populations.

\section{Resident persistence}\label{an04impv1:resident}
Initially the inactive fraction is
\(\tfrac12[1-F_b(H(\zeta))]\), where
\(F_b(l)=\frac4\pi\int_0^{l/b}(1+u^2)^{-2}\dd u\) for \(l>0\).
This is positive for every finite \(b>0,\zeta>a\).

\begin{theorem}\label{an04impv1:mass}
For each such parameter pair there is \(C_*>0\), independent of
\(T,q\), with
\[
 \mu_0\{\text{labels with }x(s)\ge\zeta\text{ for every }0\le s\le T\}
 \ \ge\ \frac{C_*}{1+T}.
\]
In particular at \(T=(2/\lambda)\ell\) the guaranteed persistent mass
is at least \(C_* /[1+(2/\lambda)\ell]\). This is an estimate in the
defined zero-weak-mass system on that time scale; it does not locate
a weak switch.
\end{theorem}
\begin{proof}
First pay the donor costs with no persistence assumption. Put
\(L=\lambda/2\) and \(w_f=(\sqrt{4+f^2}-f)/2\).
The GOE upper comparison, applied to the rank with at least upper mass
\(f\), gives
\[
 2D^+\binom{x_+}{y_+}\le
 \lambda\begin{pmatrix}-f&1\\1&0\end{pmatrix}\binom{x_+}{y_+}
 +\rho\phi(0)\binom11 .
\]
Indeed \(S(x)\ge f K(\rho x)\), \(Y\ge0\), and the nonnegative-face
test holds for the positive parts. Multiply by \((w_f,1)\).
Since \(1/2<w_f\le1\) and \(L w_f\ge L/2\), scalar integration gives
\[
 x_+(s,f)+y_+(s,f)\le C(1+\xi(f))e^{Lw_fs},\qquad
 Lw_f\le L-\lambda f/8.
\]
The constant is uniform in rank, including the bottom plateau.
The exponent bound follows by \(\sqrt{4+f^2}\le2+f^2/4\) for
\(0\le f\le1\). From the GOE density and bounded lower support,
\(\xi(f)\le C_b(1+f^{-1/3})\). Consequently the full donor moments obey
\begin{equation}
 \max(\rho K_s(s),Y(s))<B e^{Ls}(1+s)^{-2/3}
 \label{an04impv1:budget}
\end{equation}
for a fixed \(B\). To see this, integrate
\((1+f^{-1/3})e^{-\lambda fs/8}\) over \(0<f<1\); for \(s\ge1\)
the substitution \(u=fs\) bounds it by \(Cs^{-2/3}\), and the
integral is bounded on \(s\le1\). Absorb the constant in
\(K_s\le\phi(0)+\rho\int x_+\) by increasing \(B\).
This use of the ceiling includes the outliers in their own donor moments.

Fix \(T\), and set
\[
 R_T=\zeta+1+3B[(1+T)^{1/3}-1],\qquad
 w(s)=e^{Ls}\{R_T-3B[(1+s)^{1/3}-1]\}.
\]
Consider the upper ranks initially at \(x=y\ge R_T\). On
\(0\le s\le T\) one has \(w\ge e^{Ls}(\zeta+1)\) and
\(2w'=\lambda w-2B e^{Ls}(1+s)^{-2/3}\).
On the face \(x=w,\ y\ge w\),
\(2U\ge\rho K(\rho w)-\rho K_s\); on the face \(y=w,\ x\ge w\),
\(2V_s\ge\rho K(\rho w)-Y\). Since \(\rho K(\rho w)>\lambda w\),
\eqref{an04impv1:budget} makes both faces strictly inward relative
to \(w'\), including their corner. Entry and terminal faces give
the claimed finite-time invariant region. The tail asymptotic implies
\(\mu_0\{x\ge R\}\ge c_b/(1+R)^3\) for \(R\ge\zeta+1\) by continuity
and positivity on compact radii. Since \(1+R_T\le C(1+T)^{1/3}\)
the result follows. The region may even be charged with the initial
band \([R_T,2R_T]\) after decreasing the constant, by the same tail
asymptotic. The protected ranks vary with \(T\), consistently with GOE gap.
\end{proof}

\paragraph{Initial return and the uniform-in-time question.}
GOE threshold gives outward gate velocity for \(b<b_c(\rho,\zeta)\),
zero at equality, inward for \(b>b_c\). The continuous Q1 density is
positive at the initial gate. Continuity in label and time thus gives
a positive mass of initially inactive labels returning inside at short
time for \(b>b_c\) (and an initially active band crossing outward
for \(b<b_c\)). Equality has no first-order decision. Complete
finite-time return of all initially inactive labels is excluded for
every \(b\) by Theorem~\ref{an04impv1:mass}.

For a stronger fixed-fraction result the proposed wedge \(x\ge\zeta,\
y\ge x\) requires at the gate \(2U\ge0\) and at \(y=x\),
\[
 2(V_s-U)=\rho S(x)-Y\ge0 .
\]
At aligned entry \(S(x)\le K_s,\ \rho K_s<Y=m\) since the law
has positive mass above \(a\) and \(\rho K(\rho r)<r\) there. The
second face points out for all \(b\), even for very high labels.
The expanding-orthant proof instead spends a normalized donor budget
of order \(\int_0^T(1+s)^{-2/3}ds\); its upper estimate diverges
and does not certify a fixed \(R_T\) for all time.

For the alternative output-speed route let \(G=\rho K(\rho x)-Y\)
on a candidate fixed upper rank. At \(G=0\),
\[
 2G'=2\lambda\Phi(\rho x)U+Y-\rho K_s .
\]
The needed inward sign is \(2\lambda\Phi(\rho x)U+Y\ge\rho K_s\)
at every such contact. Neither \(Y>0\), order, nor the positive
output-gap estimate supplies it. In addition, keeping \(G\ge0\)
would still require \(U(\zeta,y)\ge0\) to protect the input gate.
We stop this route at those inequalities. There is no classification
here of a time-independent positive surviving fraction, or of eventual
return of each individually fixed rank. New outward crossings in RPS
and initial inward crossings are compatible with the lower bound.

\section{Bottom ranks}\label{an04impv1:bottom}
At entry \(2y'=\rho K(\rho x)-m\); it changes from negative to
positive at the unique root above \(a\), since \(m>a\).
For the input, \(2x'=\rho[K(\rho x)-S_b(x)]\) on the aligned line.
It is strictly negative at \(a\), positive for all sufficiently
large \(x\), and has a unique upward zero by the offset derivative
test in GOE threshold. At the output-speed zero the input speed is
already positive because \(\rho S_b< m\). Thus the bottom atom and
an adjacent band initially move down in both coordinates for every
\(b>0\). By continuity the direction persists on some positive
parameter-dependent interval.

There is an exact simplification at the bottom rank. Write its
coordinates \(x_a,y_a\). Since every donor has \(\tilde x\ge x_a\),
\(S(x_a)=\phi(\rho x_a)+\rho X_m\Phi(\rho x_a)\), whence
\begin{equation}
 2x_a'=\lambda\Phi(\rho x_a)(y_a-X_m),\qquad
 2y_a'=\rho K(\rho x_a)-Y.
 \label{an04impv1:bottomfield}
\end{equation}
At zero time the two velocities doubled are
\(-\lambda\Phi(\rho a)(m-a)\) and \(a-m\). In fact a crossing
of \(x=0\) is provable for large scales. Here is one explicit,
nonoptimal sufficient range. Let \(C=100e^7,\ h=1/(4C)\). Require
\[
 m>\max(1,4a,8a/(\lambda h)).
\]
This holds eventually because \(m/b\to1/\pi\) by GOE. For
\(0\le s\le1\), \(D(s)\le(1+2m)e^7\). The full field then gives
\(|X_m'|,|Y'|\le6D(s)\le Cm\). Hence through \(h\) both signed
means stay at least \(3m/4\). At \(x_a=a,y_a\le a\) the first
equation points down, and at \(y_a=a,x_a\le a\) the second does
also. The starting corner is strictly inward to the lower rectangle.
There is no finite-time escape through infinity by the growth bounds.
While \(0\le x_a\le a\), \(x_a'\le-\lambda m/8\);
throughout, \(y_a'\le-m/4\). Thus the Q4 atom crosses \(x=0\)
strictly before \(h\), and the sign at that face keeps it negative
through \(h\). A positive adjacent Q1 band also has negative input
at \(h\) by continuity in the initial coordinates in this field.
For other scales the signed bottom field and initial descent do
not decide a later crossing: control of \(y_a-X_m\) and
\(\rho K(\rho x_a)-Y\) is still needed beyond their proved range.
The growing output differences alone give no absolute bottom sign.

\section{Weak-coupled gate and ceiling}\label{an04impv1:weak}
In this section \(V=\int v\,d\mu_w\) is the signed weak output
moment and \(K_w=\int K(u)d\mu_w\ge0\). For the full WIT leading
state the strong receiver velocities are
\begin{align*}
 2U&=-(1-M)x+\rho\phi(\rho x)-\rho S(x)
                  +\lambda[V+(1-N)y]\Phi(\rho x),\\
 2V_s&=-(1-M)y-Y-K_w+\rho(1-N)K(\rho x).
\end{align*}
The weak overlap \(\mathcal W(u)=\int F(u,\tilde u)d\mu_w\)
enters weak receiver dynamics; no such term appears directly in
these strong equations. In particular at the aligned input gate
\(x=y=\zeta\) with strong entry shape \(\langle\cdot\rangle_b\)
normalized to mass \(M>0\), define \(P_\zeta=\Phi(\rho\zeta)\) and
\[
 D_{\rm gate}=\rho K(\rho\zeta)-(1-M)\zeta+
                    \lambda(V-N\zeta)P_\zeta .
\]
Then \(2U=D_{\rm gate}-\rho M S_b(\zeta)\). The increment relative
to zero weak mass is outward iff \(V>N\zeta\) (on a coherent weak
atom, iff \(v>\zeta\)). The total direction has a different test:
\begin{itemize}
\item If \(D_{\rm gate}\le Ma\), the gate velocity is inward for
every \(b>0\).
\item If \(D_{\rm gate}>Ma\), there is exactly one
\(b_*(M,N,V,\rho,\zeta)>0\) with
\(S_{b_*}(\zeta)=D_{\rm gate}/(\rho M)\).
The gate is outward below, stationary at, and inward above \(b_*\).
Equivalently at fixed \(b,M,N\) the outward half-space is
\[
 V>N\zeta+\frac{\rho M S_b(\zeta)+(1-M)\zeta-
                         \rho K(\rho\zeta)}{\lambda P_\zeta}.
\]
\end{itemize}
Here GOE gives continuity and strict increase from \(a/\rho\) to
infinity for \(S_b\). Both the explicit \(-(1-M)\zeta\) term and
the mass scaling of the donor are necessary. At a regular gate,
the active-rank velocity is \(-U/x_\sigma\); these surfaces concern
the initial direction and the corresponding initial rank flux.

\begin{lemma}[Forced per-rank ceiling]\label{an04impv1:ceiling}
On any finite chart-valid leading interval with \(M>0,\ 0\le N\le1\),
nonnegative measures and integrable moment charges, let a strong
receiver have normalized upper fraction \(f\) in a co-ordered strong
family. Define \(z=(x_+,y_+)^T\). Then
\begin{align}
 2D^+ z &\le
 \left[-(1-M)I+\lambda
 \begin{pmatrix}-Mf&1-N\\1-N&0\end{pmatrix}\right]z+d_w,
 \label{an04impv1:matrix}\\
 d_w&=\binom{\rho\phi(0)+\lambda V_+}
             {\rho(1-N)\phi(0)+(Y+K_w)_-}.\notag
\end{align}
Put \(r=-(1-M)/2+(\lambda/4)
[\sqrt{(Mf)^2+4(1-N)^2}-Mf]\). For \(s\ge t\) in that interval,
\[
 |z(s)|_2\le e^{\int_t^s r}|z(t)|_2+
      \tfrac12\int_t^s e^{\int_u^s r}|d_w(u)|_2\,du .
\]
\end{lemma}
\begin{proof}
The normalized strong weights are unchanged by the common reaction.
For \(0\le N\le1\) the two off-diagonal receiver derivatives are
\(\lambda(1-N)\Phi(\rho x)/2\ge0\), preserving co-order on this
interval. Hence \(S(x)\ge Mf K(\rho x)\).
For \(x\ge0\) use \(K(\rho x)\ge\rho x\); for the output use
\(K(\rho x)\le\phi(0)+\rho x_+\).
The factor \(\Phi\) allows a positive upper charge \(\lambda V_+\)
for both signs of \(V\). Simply replacing \(V\Phi\) by \(V\) in
an upper bound requires \(V\ge0\). When \(Y+K_w\ge0\) the bottom
forcing charge vanishes, but nonnegative \(Y\) is no longer automatic:
\[
 2Y'=(1-2M)Y-MK_w+\rho(1-N)K_s .
\]
This identity or a separate moment bound is needed to pay a
negative-part charge over time. On a zero positive-part face the
right side of \eqref{an04impv1:matrix} is nonnegative; negative
coordinates cause no uncharged term. Take the Euclidean norm
inequality using the top eigenvalue of the symmetric matrix,
including the forcing bound at \(z=0\), and integrate. This avoids
differentiating a time-dependent eigenvector.
\end{proof}
At \(M=1,\ N=0,\ V=K_w=0,\ Y\ge0\), the rate reduces to GOE's
\(r_+(f)\). At \(N=1\) the off-diagonal strengths vanish. The
forced estimate retains its time integral, even if its homogeneous
rate is zero. When \(N>1\) the displayed positive-part comparison
is not justified by that argument. On the upper faces of the resident
barrier a weak response changes both the input coefficient and the
output donor charge; the zero-weak-mass estimate cannot be continued
through a weak passage by using the sign of \(V-N\zeta\) at one time.

\section{End state and handoff}\label{an04impv1:handoff}
For any fixed \(\tau>0\) in the \(N=0\) continuation, at
\(T=\tau\ell\) the strong measure still has its co-ordered rank law,
the original atom's rank weight \(1/2\), and an initially inactive
upper mass at least \(C_* /(1+\tau\ell)\) which has stayed inactive.
The estimate is strongest as a statement about the defined system:
a diverging scaled upper radius can exceed the range of a uniform
finite-noise chart approximation. No fixed positive terminal inactive
mass, active fraction, terminal atom location, or actual switch time
is identified by these bounds.

S7d needs the strong rank profile, donor moments, active exposure and
signed lag at a matched weak switch; a weak seed and joint weak law
with their clock and field errors; and evolution of \(S,Y,K_s,V,K_w,
\mathcal W\) with rate-weighted producers through the coupled passage.
The candidate \(N_0\asymp S_0^{1-\lambda}\asymp q^2\) only motivates
the comparison time. Matching of escaped angles to an outer chart,
trained entry and radial weights including other sectors, control
of probe output and rate through saturation, and lag-relative
precision remain separate transfer obligations.
\end{document}
